\section{Method}
\label{sec:method}

\subsection{Problem Setting and $2\times2$ Factorial Design}
\label{sec:method:design}

%% [§3.1 ¶1]
For a multi-hop question $q$ from dataset $d$, the agent uses an LLM $M$ and corpus
$\mathcal{C}_d$. Its state $s=(q,\,Q_{\mathrm{sub}},\,E)$ tracks the question,
sub-questions, and collected evidence. The \emph{evidence source}
$e\in\{\cond{recall},\cond{rag}\}$ determines where evidence comes from. The \emph{loop
controller} $c\in\{\cond{llm},\cond{laya}\}$ decides when to stop and answer. Their
combinations give the four conditions in
Table~\ref{tab:conditions}.

%% [§3.1 ¶2]
\begin{table}[t]
  \centering
  \caption{The four conditions of the $2\times2$ factorial design. All conditions share
  the loop of Algorithm~\ref{alg:loop}.}
    \label{tab:conditions}
  \small
  \begin{tabular}{|l|l|l|}
    \hline
    Condition & Evidence source $e$ & Controller $c$ \\
    \hline
    Recall+LLM  & \cond{recall} (closed-book) & \cond{llm} \\
    \hline
    RAG+LLM     & \cond{rag} (retrieve + re-rank) & \cond{llm} \\
    \hline
    Recall+Laya & \cond{recall} (closed-book) & \cond{laya} \\
    \hline
    RAG+Laya    & \cond{rag} (retrieve + re-rank) & \cond{laya} \\
    \hline
  \end{tabular}
\end{table}

%% [§3.1 ¶3]
The design separates the two main effects and their interaction. Recall+LLM also uses
the loop: the LLM asks sub-questions, answers from memory, and decides when to stop.
It does not give a single-shot closed-book answer. Differences between \cond{llm} and
\cond{laya} therefore reflect who stops the loop and which tools are available, rather
than iterative decomposition.

\subsection{Shared Decomposition Loop}
\label{sec:method:loop}

%% [§3.2 ¶1]
Algorithm~\ref{alg:loop} shows the shared loop. At each hop, the controller can stop.
Otherwise, the LLM asks a sub-question and collects evidence through retrieval or recall.
The LLM answers from this evidence when the controller stops or the loop reaches
$H_{\max}=4$ hops.

%% [§3.2 ¶2]
\begin{algorithm}[t]
  \caption{Decomposition loop shared by the four conditions. Retrieved passages are
  de-duplicated by passage id and capped at $|E|\le8$.}
    \label{alg:loop}
    \raggedright% the wrapped Require line without stretched spaces

  \begin{algorithmic}[1]
    \Require question $q$; evidence source $e\in\{\cond{recall},\cond{rag}\}$;
             controller $c\in\{\cond{llm},\cond{laya}\}$; $H_{\max}=4$
    \State $Q_{\mathrm{sub}}\gets[\,]$; \ $E\gets[\,]$
    \For{$h=1,\dots,H_{\max}$}
      \State $d_h\gets\varnothing$ \Comment{no hint unless $c=\cond{laya}$}
      \If{$c=\cond{laya}$}
        \State $d_h\gets\mathrm{Laya}(\mathrm{BuildState}(q,Q_{\mathrm{sub}},E))$
        \IfThen{$d_h.p_{\mathrm{suf}}\ge0.5$}{\textbf{break}} \Comment{enforced}
      \EndIf
      \State $q_h\gets M.\mathrm{SubQuestion}(q,Q_{\mathrm{sub}},E,d_h)$
      \If{$c=\cond{llm}$ \textbf{and} $M$ called \texttt{final\_answer}}
        \State \textbf{break}
      \EndIf
      \State append $q_h$ to $Q_{\mathrm{sub}}$
      \If{$e=\cond{rag}$}
        \State $E\gets E\cup\mathrm{Rerank}(q_h,\mathrm{Dense}(q_h,20))[1{:}3]$
      \Else
        \State $E\gets E\cup\{M.\mathrm{Recall}(q_h)\}$
      \EndIf
    \EndFor
    \State \Return $M.\mathrm{Answer}(q,E)$
  \end{algorithmic}
\end{algorithm}

%% [§3.2 ¶3]
\textbf{Controller.} When $c=\cond{laya}$, Laya runs at each hop's start. If its
sufficiency probability is $p_{\mathrm{suf}}\ge0.5$, the loop ends and the LLM must
answer. Otherwise, its predicted next action and remaining hops guide the LLM, with
only \texttt{ask\_sub\_question} available. When $c=\cond{llm}$, the LLM can call
\texttt{ask\_sub\_question} or stop through \texttt{final\_answer}. We count the
stopping hop, so hops equal decisions. Each
question has a 300\,s budget.

%% [§3.2 ¶4]
\textbf{Evidence and failures.} When $e=\cond{rag}$, the dense retriever uses the
sub-question to find the top-20 passages in $\mathcal{C}_d$. The re-ranker's top-3
passages are added to $E$. When $e=\cond{recall}$, we add the LLM's closed-book answer
to the sub-question. A reply without a valid call to an available tool, including a
plain-text reply, ends the question as \texttt{invalid\_tool}. Exceeding the time budget
gives \texttt{timeout}. Other failures receive \texttt{error} or \texttt{oom}. All count
as incorrect.

\subsection{Laya as the Decision Layer}
\label{sec:method:laya}

%% [§3.3 ¶1]
At each hop, Laya answers three fixed English questions about the state: whether to stop,
what to do next, and how much evidence is missing. Each uses a different question type.

%% [§3.3 ¶2]
\begin{itemize}
  \item \texttt{sufficient} (\texttt{noul}): is the evidence sufficient for the original question? Laya returns $p_{\mathrm{suf}}=P(\text{true})$.
  \item \texttt{next\_action} (\texttt{choice}): (A) answer now; (B) form one more
        sub-question and search; (C) rephrase the query and search again.
  \item \texttt{remaining\_hops} (\texttt{score}): how many key evidence passages are
        missing? The options are 0, 1, or 2 (two or more).
\end{itemize}

%% [§3.3 ¶3]
\textbf{State construction.} \textsc{BuildState} converts the question, sub-questions,
and evidence into text within a 900-token limit, keeping only the first two sentences of
each item. It uses the same format for training and inference.

\subsection{Weak Supervision by State Simulation}
\label{sec:method:weak}

%% [§3.4 ¶1]
The datasets lack human labels for these three decisions, but their training splits
identify the supporting paragraphs. For a training question with gold paragraphs $G$ and
distractors $D$, we simulate four kinds of states: (1) \emph{question only}, $E=\emptyset$;
(2) \emph{partial gold}, a non-empty proper subset of $G$ plus 0--2 passages from $D$;
(3) \emph{all gold}, all of $G$ plus 0--2 passages from $D$, shuffled; and
(4) \emph{distractors only}.

%% [§3.4 ¶2]
Let $m=|G\setminus E|$ be the number of gold paragraphs missing from the evidence. The
labels are
\begin{align}
  \texttt{sufficient}      &= \mathbf{1}[m=0], \\
  \texttt{next\_action}    &=
    \begin{cases}
      \text{A} & \text{if } m=0,\\
      \text{C} & \text{if } E\neq\emptyset \text{ and } E\cap G=\emptyset,\\
      \text{B} & \text{otherwise},
    \end{cases} \\
  \texttt{remaining\_hops} &= \min(m,\,2).
\end{align}
Missing gold evidence calls for another decomposition step. If the evidence contains only
distractors, the query should be rephrased instead.

%% [§3.4 ¶3]
With a fixed seed, we sample 800, 100, and 100 training questions per dataset for
decision-layer training, calibration, and held-out evaluation. We target about
30{,}000 training decisions. Before the main runs, \texttt{sufficient} and
\texttt{next\_action} must each reach 0.75 accuracy and exceed the zero-shot base
checkpoint (evaluated only here) by at least 0.20. Mean absolute error (MAE) on
\texttt{remaining\_hops} must be at most 0.5.

\subsection{Evaluation Protocol}
\label{sec:method:eval}

%% [§3.5 ¶1]
Our primary measure is EM after official answer normalisation. We also report token-level
F1 (for MuSiQue, the maximum over the answer and its aliases). For each question, we record
end-to-end latency, device-wide peak GPU memory, LLM prompt and output tokens, hops, and
Laya-call latency. For the decision layer, we report accuracy, Brier score, ECE, and
latency per decision. EM differences include paired percentile-bootstrap 95\% CIs based
on 10{,}000 resamples of the 200 questions.

\subsection{Implementation Details}
\label{sec:method:impl}

%% [§3.6 ¶1]
\textbf{Models and Datasets:} We use five instruction-tuned Qwen3.5
models~\cite{qwen35} (0.8B, 2B, 4B, 9B, and 27B parameters) with thinking disabled,
an 8192-token context, greedy decoding, and seed 42. All run on one 24\,GB consumer GPU
(software, checkpoints, and licenses are in Appendix~\ref{sec:appendix}). We evaluate
HotpotQA~\cite{yang2018hotpotqa} in the distractor setting,
2WikiMultihopQA~\cite{ho2020twowiki} (2Wiki), and the answerable subset of
MuSiQue~\cite{trivedi2022musique}. For each dataset, we sample 200 test questions from the
validation split and the 800/100/100 decision-data questions from the training split.
No test questions enter decision-layer training, calibration, or held-out evaluation.
Before the main runs, each model answers 20 fixed tool-calling probe items. We retain
all models.

%% [§3.6 ¶2]
\textbf{Retrieval:} Following~\cite{yang2025enhancing}, we embed passages and
sub-questions with BGE-M3~\cite{chen2024bgem3}, retrieve passages by exact inner-product
search, and re-score them with the BGE cross-encoder re-ranker~\cite{bgerankerv2m3}.
Each dataset's corpus combines all supporting and distractor paragraphs from its 200 test
questions. We remove duplicates and use one chunk per paragraph. The embedder, re-ranker,
and Laya share the GPU with the generator.

%% [§3.6 ¶3]
\textbf{Controller:} We fine-tune one shared decision layer for all datasets from the
multilingual Laya checkpoint~\cite{laya}. Training uses Laya's reinforcement learning
with strictly proper scoring rules~\cite{laya}. We then fit one temperature per question
type on the calibration split (details are in Appendix~\ref{sec:appendix}).

